\documentclass[11pt]{article}

\usepackage{maiacustom}

\begin{document}

\psettitle{Astronomy Question Bank}

These questions were carefully written or selected to prepare students for the astronomy team selection process in Brazil. Some questions are not original; those are marked with () before the statement. The question-bank template is the same one used by Professor Kevin Zhou. His work is invaluable, and many ideas in this set can be found in his handouts.

% \begin{psidea}{Idea title}{}
% Idea
% \end{psidea}

% \begin{psexample}{Example title}{}
% Example
% \end{psexample}

% \begin{pssolution*}{}{}
% Solution
% \end{pssolution*}

% \begin{psremark*}{Remark title}{}
% Remark
% \end{psremark*}
\section{Cosmology}

\pts{5}
\begin{pproblem}
    The Friedmann equation is one of the most important equations for the study of cosmology and of the Universe. The goal of this problem is to derive the fundamental equations of cosmology. First, let us begin with some definitions:
    \begin{enumerate}[label=\roman*)]
        \item Because of the expansion of the Universe, the distance between two points is given by
        \[r(t) = r_0 a(t)\]
        where \(a(t)\) is known as the scale factor and \(r_0\) is the distance measured at \(t = 0\). 
        
        \item The Universe follows the Robertson-Walker metric, which can be simplified to
        \[-c^2 dt^2 + dr^2 = 0\]
        \begin{alternativas}
            \item Find an expression for the comoving distance, \(r_0\), in integral form, from the definitions given above.
            \\
            \item The first Friedmann equation has the form
            \[H(t)^2 \equiv \left(\frac{\dot{a}}{a}\right)^2 = k_1 \varepsilon(t) + \frac{C}{r_0^2 a^2}\]
            Find the value of the constant \(k_1\) for a spherical universe, in terms of fundamental constants. In the equation above, \(\varepsilon(t)\) is the energy density of the universe, \(C\) is a constant related to its energy, and \(a\) is the scale factor.

            \textbf{Hint:} You can obtain this equation either from energy conservation or by using Newton's second law.

            \item Repeat the previous item for a cylindrical universe expanding radially. The equation you find should have the form

            \[\left(\frac{\dot{a}}{a}\right)^2 = f(a) \varepsilon(t) + \frac{C'}{r_0^2 a^2}\]

            Find the function \(f(a)\).

            \item Returning now to a ``normal universe''. In cosmology, we define the critical energy density as the energy density of a universe in which \(C = 0\). Find an expression for the critical density.
            \item Rewrite the first Friedmann equation in terms of
            \[\Omega(t) = \frac{\varepsilon(t)}{\varepsilon_c(t)}\]
            where \(\varepsilon_c(t)\) is the critical energy density.
        \end{alternativas}
    \end{enumerate}
\end{pproblem}

\pts{3}
\begin{pproblem}(Adapted from NAO 2019) 
    Consider a flat universe in which the gravitational constant is no longer constant and is instead defined by
    \[G(a) = G_0 f(a)\]
    where \(f(a)\) is a function of the scale factor.
    \begin{alternativas}
        \item What would the Friedmann equation be in this universe? Assume that the universe is flat, \(C = 0\), and that it is composed only of baryonic (``bright'') matter. Leave your answer in terms of \(H_0\), \(f(a)\), \(a\), and \(\Omega_{m,0}\), where \(H_0\) is the value of the Hubble constant at the present time, and \(\Omega_{m,0}\) is the density parameter,
        
        In the case \(f(a) = e^{b(a-1)}\), where \(b=2\).
        
        \item Estimate the age of this universe, assuming that it consists only of baryonic matter (``bright'' matter).
                
        \item What is the behavior of the age when \(t \rightarrow \infty\)?
    \end{alternativas}

    You may find the following relations useful:
    \begin{paracol}{2}
        \[\int_0^\infty x^2 e^{-x^2} = \frac{\sqrt{\pi}}{4}\]

        \switchcolumn

        \[\int_0^1 x^2 e^{-x^2} \approx 0.189\]
    \end{paracol}
\end{pproblem}

\pts{5}
\begin{pproblem}(Problem set 8 - 2021)
    The Friedmann equation is given by
    \begin{equation}
        H^2 = \left( \frac{\dot{a}}{a} \right)^2 = \frac{8\pi G \varepsilon}{3c^2} - \frac{k c^2}{a^2},
    \end{equation}

    where \(a\) is the scale factor at time \(t\), \(\varepsilon\) is the energy density at time \(t\), and \(k\) is the parameter that characterizes the geometry of the universe and may take any real value. Considering a universe composed only of non-relativistic baryonic matter and solving this nonlinear differential equation for \(k > 0\), one obtains the following solutions in terms of the parameter \(\theta \in [0, 2\pi]\):

    \begin{align}
        a(\theta) &= \frac{4\pi G \varepsilon_0}{3k c^4} \left( 1 - \cos\theta \right), \\
        t(\theta) &= \frac{4\pi G \varepsilon_0}{3k^{3/2}c^5} \left( \theta - \sin\theta \right).
    \end{align}

    Consider a universe with \(\Omega_0 = 4\) and \(H_0 = 67.4 \ \text{km/s/Mpc}\).

    \begin{alternativas}
        \item Starting from the Friedmann equation, show that \(k c^2 = H^2 \left( 1 - \Omega \right) a^2\). Finally, rewrite the parametric equations for \(a\) and \(t\) in terms of the present density parameter \(\Omega_0\), the present Hubble constant \(H_0\), and the parameter \(\theta\). Do not substitute their respective numerical values.
        
        \item Find the age \(t_0\) of this universe in terms of \(H_0\), and then in billions of years.
        
        \item The so-called \textit{lookback time}, \(\Delta t_L\), is how far in the past the universe had a given scale factor \(a\). What is \(\Delta t_L\), in billions of years, for when the size of the universe was \(1/3\) of what it is at present?
        
        \item Determine \(\theta_n\), and then \(t_n\), for which \(H = 0\).
    \end{alternativas}
\end{pproblem}
\end{document}
